\section{ダイナミクス}
\label{sec:dynamics}

\subsection{沈殿}
沈殿とは、反復的なトークンと場の相互作用のもとで井戸が深まり、概念が安定化するという提案的過程の呼称である。
以下の三段階は機能的なフィードバック構成を記述するものであり、実行された動的過程の実証ではない。
最初のトークン・盆地の立ち上がりと、結合の更新上の帰結は未解決である。

\textbf{段階1：初期摂動。}
提案するシナリオでは、新たな入力$\Delta$を、$\sigma_{\mathrm{drive}}$のスパイク、$I$の局所的増加、一過性のポテンシャル勾配と関連づける。
情報流$\JI$は摂動に応答するものとしてモデル化するが、これは再構成された思考経路ではない。

\textbf{段階2：強化。}
類似入力へのその後の接触や、内的な再帰的活性化によって摂動が繰り返されると、$\Phi$の局所勾配が強化されると仮定する。
井戸の深化はトークンと場の結合\eqref{eq:coupling}の提案的効果であり、個々の再利用や表式だけによって保証される結果ではない。

\textbf{段階3：安定化。}
提案する安定化段階では、井戸の深化を通過コストの低下と、その後の再利用をさらなる強化と関連づける。
自己維持する概念は、このように得られる維持仮説であり、フィードバック構成から実証された帰結ではない。

\subsection{忘却}
ここでは忘却を、再利用頻度が低下すると$\Phi$がより平坦な配置へ緩和することとの類比で解釈する。
既存の散逸項を形式的テンプレートとして保持する。
\begin{equation}
  \sigma_{\mathrm{dissip}} = \Gamma C.
\end{equation}
このモデル描像では、アクセス可能性の低下を、保存された記号の消去ではなく、安定性の再分配と関連づける。
浅くなる井戸では通過コストが高まると提案する。
これは、忘れられた概念が決して消滅しないことを確立しない。
再活性化（想起）も、再開されたトークンと場の相互作用を通じた再深化としてモデル化するが、普遍的な記憶機構とは同一視しない。

\subsection{概念変化}
確立された概念の意味、範囲、適用の変化である概念変化を、ここでは$\Phistar$の変化を用いて表現する。
二つの井戸の融合、一つの井戸の二つへの分裂、地形内での変位などである。
これらの幾何的選択肢は認知的対応の提案を整理するものであり、長時間対象を確定したり、トポロジーから認知への写像を証明したりするものではない。

この候補となる記述では、そのような変化を既存アトラクタの深さに対する持続的摂動と関連づけるが、信念改訂の較正された必要条件や十分条件とはしない。
有効障壁項$\sigma_{\mathrm{barrier}}$は、この有効記述において保護的な役割を担う。
$C\to\Cmax$となるにつれて、抵抗という解釈を、深く定着した概念を変える困難さの増大と比較する。
これは概念的観測窓内の過閉包に対する有効記述として理解すべきであり、外部から課された壁ではない。
この記述は信念改訂への抵抗についての候補となる説明を与えるが、逆数項は認知的な禁止を確立しない。

\subsection{洞察と急速遷移仮説}
問題が解けた、あるいは離れていた考えが結びついたという突然の感覚は、$\Phi$の局所トポロジーの急速な変化との対応の提案を動機づける。
井戸の融合や障壁の崩れは、分岐に似た表現を与える。
幾何的遷移、認知的な問題解決、感じられる出来事は別々の主張であり、ここでは分岐定理も独立した認知的終点も与えない。

突然性の主観的感覚には、摂動の蓄積に続いて急速な遷移が起こるという閾値描像との対応を仮定する。
この描像は、知覚的閾値を独立に同定するものでも、遷移と名づけるだけで感じられる突然性を説明するものでもない。
